\documentclass[12pt]{article}

\usepackage[utf8]{inputenc}
\usepackage[T1]{fontenc}
\usepackage{lmodern}
\usepackage{amsmath}
\usepackage{booktabs}
\usepackage{graphicx}
\usepackage{geometry}
\usepackage[hidelinks]{hyperref}
\usepackage{natbib}
\usepackage{longtable}

\geometry{margin=1in}
\graphicspath{{Paper/figures/}{figures/}}
\newcommand{\tablepath}{Paper/tables/}
\newcommand{\figurelang}{en}
\newif\ifpolicyreport
\policyreportfalse

\title{Beyond Formality: The Firm-Size Wage Premium in Colombia}
\author{}
\date{}

\begin{document}

\maketitle

\begin{abstract}
Firm-size wage premia are well documented in advanced economies, but their interpretation in middle-income labor markets with large informal sectors remains less settled. This paper documents the firm-size wage premium in Colombia using household survey data for 2008 to 2025. Employment is concentrated in solo work and micro-firms, while real hourly wages rise sharply with firm size. After controlling for gender, age, education, formality, and sector-year fixed effects, workers in firms with more than 100 employees earn about 43 percent more than comparable solo workers. The premium is positive throughout the sample period, with no evidence of secular decline, and is not simply a relabeling of the formal-informal wage gap. The size gradient is steep within informal employment, while the large-firm advantage within formal employment is narrower. These patterns are consistent with frictions that may constrain firm growth and worker reallocation.
\end{abstract}

\vspace{0.5em}
\noindent\textbf{Keywords:} Firm-size wage premium; informality; Colombia.

\vspace{0.5em}
\noindent\textbf{JEL Codes:} J31; J46; O17; L25.


\section{Introduction}

Firm size is a well-established correlate of wages in labor economics \citep{BrownMedoff1989,OiIdson1999,troske1999evidence}. In labor markets where many workers are own-account workers or employed in micro-firms, wage differences between small and large firms help identify where questions about firm growth, worker allocation, and wage inequality are most relevant. This paper asks whether Colombian workers in larger firms earn higher real hourly wages than observably similar workers in smaller firms, how much of that gap is explained by observable worker characteristics, whether the gap has increased or decreased over time, and whether the wage gap between small and large firms differs between formal and informal employment.

Our motivation to measure the firm-size wage premium in Colombia is both empirical and policy-oriented. A recent policy diagnosis argues that low-income households are disproportionately linked to own-account work, very small businesses, informality, and low-earning jobs, whereas higher-income households are more often connected to larger and more formal employers \citep{eslava2026crecimiento}. This diagnosis is suggestive, but it does not show whether firm size remains associated with higher wages after accounting for worker characteristics and formality. The observed link between larger employers and higher household income may instead reflect the fact that larger firms employ workers with different characteristics, or that larger firms are more likely to offer formal jobs. This distinction matters because the wage gap associated with firm size has not been systematically documented in Colombia. The question is especially relevant in Colombia because employment remains concentrated outside the largest-firm categories, while recent growth theory links the rise of large firms to economic development \citep{Chen2026}.

We study the firm-size wage premium in Colombia using harmonized microdata from the Gran Encuesta Integrada de Hogares (GEIH) for 2008--2019 and 2021--2025. We construct a comparable firm-size classification, measure real hourly labor income, document the distribution of employment across firm sizes, and estimate wage gaps that progressively condition on sector-year effects, worker characteristics, and formality. Because unobserved worker characteristics may be correlated with firm size, the estimates are conditional wage gaps rather than causal effects of firm size. This interpretation follows the broader firm-size wage premium literature: a premium that survives observable controls is informative about sorting, segmentation, firm-side pay differences, and possible allocative frictions, but it is not by itself proof of any one mechanism \citep{BrownMedoff1989,OiIdson1999,AbowdKramarzMargolis1999,CardHeiningKline2013,ElBadaouiStroblWalsh2010,BalkanTumen2016}.

The evidence delivers five main results. First, employment is concentrated at the bottom of the firm-size distribution: solo workers and firms with up to ten workers account for about two thirds of employment. Second, raw wage gaps are large across firm-size categories, sectors, education groups, and formality status, while the wage gap by sex is small by comparison. Third, the firm-size wage premium attenuates substantially after adding worker controls, sector-year fixed effects, and formality, but it does not disappear: in the full specification, workers in firms with 101 or more employees earn about 43 percent more than otherwise comparable solo workers. Fourth, the conditional premium is positive throughout the period; the point estimates suggest an increase between 2008 and 2025, but the change is not statistically significant. Fifth, formality status affects the firm-size wage premium. Formal workers earn more than informal workers within every firm-size category, but the conditional firm-size wage premium is much steeper among informal workers. Within formal employment, the conditional premium relative to formal solo workers is small, although workers in 101+ firms earn more than workers in intermediate-size formal firms.

The rest of the paper is organized as follows. Section~2 describes the GEIH data, the harmonization of firm-size categories, and the construction of real hourly labor income. Section~3 presents the descriptive evidence on employment shares and wage gaps by firm size, formality, education, gender, sector, and year. Section~4 introduces the econometric specifications. Section~5 reports the main regression results, the dynamic estimates, the formality interactions, and the elasticity benchmark, contrasting each result with the relevant literature. Section~6 discusses interpretation, limitations, and policy implications. Section~7 concludes.

\section{Data}

The empirical analysis uses microdata from the Gran Encuesta Integrada de Hogares (GEIH), Colombia's nationally representative labor force survey produced by DANE. The current harmonized file covers the years 2008--2019 and 2021--2025. We exclude 2020 because the COVID-19 pandemic disrupted both labor-market outcomes and survey collection, raising comparability and data-quality concerns. The analysis therefore focuses on non-pandemic survey years. The main qualitative patterns documented below are visible on both sides of the excluded year: the firm-size wage premium is positive in the pre-pandemic period and remains positive after 2020. The working sample contains 5,353,440 worker-year observations with valid real hourly labor income and positive expansion weights. The baseline regression sample contains 5,268,514 observations after requiring non-missing values for the controls used in the econometric specification.

The sample is restricted to employed workers with positive labor income, valid weekly hours, and positive survey expansion factors. Monthly labor income is converted to annual labor income and divided by annual hours, computed as weekly hours times 52. We then express hourly labor income in constant 2025 Colombian pesos using the end-of-the-year Consumer Price Index:

\begin{equation}
  w_{it} =
  \frac{12 \times MonthlyIncome_{it}}{52 \times WeeklyHours_{it}}\times
  \frac{CPI_{2025}}{CPI_t}.    
\end{equation}

Firm size is measured using the worker-reported firm-size variable in GEIH. The construction harmonizes the categories into the following ordered groups: Solo, 2-3, 4-5, 6-10, 11-19, 20-30, 31-50, 51-100, and 101+. When the newer questionnaire separates 101-200 and 201+ workers, both categories are collapsed into 101+ to maintain comparability with earlier years. Firm size is therefore reported by workers, and the data do not identify firms.

The control variables are also harmonized across survey changes. Education is grouped into comparable categories: no education, preschool, primary, lower secondary, upper secondary, and higher education. Formality is defined using pension contribution status: workers who contribute to pensions are classified as formal, while workers who do not contribute are classified as informal. This definition follows the social-protection approach commonly used in Latin American labor-market measurement, which identifies informality through lack of access to social-security benefits such as pension or health insurance \citep{WorldBankLACInformality2026}. It is also the worker-level definition used in the recent Colombian ESPE on labor and business informality \citep{OteroCortesEtAl2025}. Economic activity is harmonized across CIIU revisions into broad sectors, which allows the regressions to absorb sector-year shocks.


\section{Descriptive Statistics}

The descriptive section is organized around the empirical facts that motivate the regression analysis. We first document where workers are located in the firm-size distribution and how wages per hour vary across that distribution. We then
place the raw firm-size wage gap next to other remuneration gaps, describe worker composition by firm size, and examine whether the raw firm-size pattern appears within formality status, sex, education groups, and economic sectors.

Table~\ref{tab:descriptive-core-size-2008-2025} establishes the two core raw facts behind the paper. First, mean real hourly income rises sharply with firm size. In 2008, solo workers earned 5,388 COL 2025 pesos per hour on average, compared with 12,943 pesos among workers in firms with 101 or more employees, a ratio of 2.4 to 1. By 2025, the corresponding means were 6,534 and 16,859 pesos, a ratio of 2.6 to 1. Appendix Tables~\ref{tab:descriptive-core-size-2008-2025-101-reference} and \ref{tab:descriptive-core-size-2008-2025-11-19-reference} report the same means using alternative reference categories. Relative to 101+ workers, solo workers earned 42 percent of the large-firm mean in 2008 and 39 percent in 2025. Relative to workers in firms with 11--19 workers, 101+ workers earned 1.49 times as much in 2008 and 1.45 times as much in 2025. The all-worker mean increased from 7,752 to 10,295 pesos, equivalent to 1.7 percent annualized real growth between 2008 and 2025. Average weekly hours in the analytic sample declined from 47.8 to 43.8 over the same period, so real hourly wages grew faster than implied labor income per worker, which increased by 1.5 percent per year. Second, employment remains concentrated at the bottom of the firm-size distribution. In 2025, solo workers accounted for 35 percent of employment, and workers in firms with up to ten workers accounted for 61 percent. The middle of the distribution was thin: firms with 11--100 workers jointly accounted for 14 percent of employment. The 101+ category grew from 19 percent of employment in 2008 to 25 percent in 2025, but this increase did not eliminate the large employment mass in solo work and micro-firms.

\input{Paper/sections/descriptive_core_size_2008_2025_table}

This shift toward the largest category is relevant beyond the wage-premium
evidence. Recent growth theory emphasizes that the right tail of the firm-size
distribution tends to thicken with economic development, and that large firms
can matter for aggregate growth because they are central to idea search,
diffusion, and productivity dynamics \citep{Chen2026}. In this light,
Colombia's increase in the employment share of 101+ firms is consistent with
some movement toward a more developed firm-size structure. At the same time, the
continued concentration of employment in solo work and micro-firms shows that
this transformation remains incomplete.

Table~\ref{tab:descriptive-wage-gaps-2025} puts the raw firm-size gap in
context by comparing it with remuneration gaps across sex, formality status,
education, and economic activities in 2025. The wage gap by sex is small in
these weighted means: men earn 0.6 percent below the overall mean and women
0.8 percent above it. Other gaps are much larger. Formal workers earn 48.0
percent above the overall mean, while informal workers earn 39.1 percent below
it; this contrast is close in magnitude to the 2.6-to-1 ratio between workers
in 101+ firms and solo workers reported in
Table~\ref{tab:descriptive-core-size-2008-2025}. Education and sector
differences are larger still. Workers with higher education earn 71.6 percent
above the overall mean, while workers with preschool or less earn 57.9 percent
below it. Across activities, agriculture is 41.4 percent below the mean,
whereas information and communication is 112.8 percent above it. The raw
firm-size gap is therefore far larger than the wage gap by sex, comparable to
the formality gap, and smaller than the largest education and sector gaps.

\begin{table}[!htbp]
  \centering
  \caption{Mean real hourly labor income gaps by worker group, 2025}
  \label{tab:descriptive-wage-gaps-2025}
  \scriptsize
  \begin{tabular}{lp{0.34\textwidth}rrrr}
    \toprule
    Dimension & Group & Employment share & Mean income & Gap & Ratio \\
    \midrule
    Sex & Men & 58.8\% & 10,236 & -0.6\% & 1.00 \\
        & Women & 41.2\% & 10,379 & 0.8\% & 1.01 \\
    \midrule
    Formality & Formal & 44.9\% & 15,241 & 48.0\% & 2.43 \\
              & Informal & 55.1\% & 6,271 & -39.1\% & 1.00 \\
    \midrule
    Education & Preschool or less & 2.9\% & 4,333 & -57.9\% & 1.00 \\
             & Primary & 18.1\% & 5,693 & -44.7\% & 1.31 \\
             & Secondary & 45.5\% & 7,085 & -31.2\% & 1.63 \\
             & Higher education & 33.5\% & 17,667 & 71.6\% & 4.08 \\
    \midrule
    Activity & A - Agriculture & 13.8\% & 6,036 & -41.4\% & 1.00 \\
           & T - Household employers & 3.1\% & 6,454 & -37.3\% & 1.07 \\
           & I - Accommodation and food & 7.6\% & 6,930 & -32.7\% & 1.15 \\
           & H - Transport & 7.8\% & 7,775 & -24.5\% & 1.29 \\
           & F - Construction & 6.9\% & 8,406 & -18.3\% & 1.39 \\
           & G - Commerce & 17.0\% & 8,535 & -17.1\% & 1.41 \\
           & R/S - Arts and other services & 5.4\% & 8,765 & -14.9\% & 1.45 \\
           & C - Manufacturing & 11.0\% & 9,591 & -6.8\% & 1.59 \\
           & D/E - Utilities and water & 1.3\% & 11,295 & 9.7\% & 1.87 \\
           & B - Mining & 1.3\% & 11,441 & 11.1\% & 1.90 \\
           & L/M/N - Business services & 9.0\% & 13,459 & 30.7\% & 2.23 \\
           & Q - Health & 4.4\% & 14,859 & 44.3\% & 2.46 \\
           & P - Education & 4.2\% & 19,980 & 94.1\% & 3.31 \\
           & K - Finance and insurance & 1.9\% & 21,389 & 107.8\% & 3.54 \\
           & O - Public administration & 3.6\% & 21,508 & 108.9\% & 3.56 \\
           & J - Information and communication & 1.7\% & 21,903 & 112.8\% & 3.63 \\
    \bottomrule
  \end{tabular}
  \vspace{0.3em}
  \begin{minipage}{0.95\textwidth}
  \footnotesize
  Notes: Statistics use 2025 observations and GEIH expansion weights. Employment shares are computed within each dimension. Mean income is real hourly labor income in constant 2025 pesos. The gap is the percent difference in each group's weighted mean relative to the overall 2025 weighted mean. The ratio divides each group's mean by the lowest mean within the same dimension. Formality rows exclude pensioned occupied workers. Sector rows omit extraterritorial organizations.
  \end{minipage}
\end{table}


Table~\ref{tab:descriptive-worker-characteristics-2025} shows why the raw firm-size wage premium cannot be interpreted as a pure size effect. Larger firms employ a very different workforce. In 2025, only 8.3 percent of solo workers were formal, compared with 46.6 percent in firms with 6--10 workers and 98.0 percent in firms with 101 or more workers. Higher education also rises strongly with firm size, from 16.7 percent among solo workers to 62.9 percent in the 101+ category. The share of women varies less systematically across firm-size groups, while solo workers are older on average. These differences motivate the comparisons that follow, which ask whether the firm-size premium is visible within sex, formality, education, and sector groups.

\begin{table}[htbp]
  \centering
  \caption{Worker characteristics by firm size, 2025}
  \label{tab:descriptive-worker-characteristics-2025}
  \small
  \begin{tabular}{lrrrr}
    \toprule
    Firm size & Women & Formal & Higher education & Mean age \\
    \midrule
    Solo & 41.3\% & 8.3\% & 16.7\% & 44.5 \\
    2-3 & 34.6\% & 13.0\% & 17.0\% & 41.3 \\
    4-5 & 36.5\% & 24.2\% & 20.2\% & 37.6 \\
    6-10 & 40.5\% & 46.6\% & 29.9\% & 37.0 \\
    11-19 & 40.4\% & 64.4\% & 38.8\% & 37.2 \\
    20-30 & 42.7\% & 81.2\% & 45.4\% & 37.0 \\
    31-50 & 41.2\% & 90.4\% & 51.0\% & 37.5 \\
    51-100 & 43.7\% & 94.3\% & 53.4\% & 37.5 \\
    101+ & 45.8\% & 98.0\% & 62.9\% & 38.0 \\
    \bottomrule
  \end{tabular}
  \vspace{0.3em}
  \begin{minipage}{0.95\textwidth}
  \footnotesize
  Notes: Statistics use only 2025 observations and are weighted by GEIH expansion weights. Higher education corresponds to workers classified as superior or university educated.
  \end{minipage}
\end{table}


Figure~\ref{fig:income-formality-comparison} looks at the firm-size wage premium by formality status. The bubble size represents the percentage of formal or informal workers in each firm-size category in the corresponding year. Formal workers earn more than informal workers and are much more concentrated in large firms. The wage pattern by firm size is still visible within informal workers: among them, mean hourly income rises from 5,677 pesos for solo workers to 12,813 pesos in firms with 101 or more workers in 2025. Among formal workers, the raw profile is closer to a U-shape: formal solo workers earned 15,990 pesos per hour in 2025, formal workers in firms with 4--5 workers earned 12,036 pesos, and formal workers in firms with 101 or more workers earned 16,944 pesos. Table~\ref{tab:formal-worker-profile-2025} shows that formal solo work is a selected but heterogeneous category: 82.1 percent are own-account workers and 50.9 percent have higher education, with a sizable concentration in professional and administrative services, but 16.2 percent are domestic workers and 15.6 percent are in transport. Thus, the raw U-shape should be interpreted as evidence of worker composition within formal solo work, rather than as evidence that all formal solo workers are high-earning professionals.

\begin{figure}[!htbp]
  \centering
  \includegraphics[width=0.95\textwidth]{fig69_\figurelang.png}
  \caption{Mean real hourly income by firm size and formality status, 2008 and 2025}
  \label{fig:income-formality-comparison}
\end{figure}

\begin{table}[htbp]
  \centering
  \caption{Profile of formal workers by firm-size group, 2025}
  \label{tab:formal-worker-profile-2025}
  \small
  \begin{tabular}{lrrr}
    \toprule
    Characteristic & Formal solo & Formal 2--100 & Formal 101+ \\
    \midrule
    Share of formal workers & 6.4\% & 38.7\% & 54.9\% \\
    Mean real hourly income & 15,990 & 12,696 & 16,944 \\
    Mean age & 44.9 & 38.3 & 37.9 \\
    Women & 43.3\% & 41.9\% & 45.7\% \\
    Higher education & 50.9\% & 48.2\% & 63.4\% \\
    Own-account workers & 82.1\% & 7.7\% & 9.6\% \\
    Domestic workers & 16.2\% & 0.9\% & 0.0\% \\
    Professional and administrative services & 22.8\% & 14.6\% & 9.4\% \\
    Transport and storage & 15.6\% & 5.5\% & 5.6\% \\
    Commerce and repair & 14.7\% & 21.1\% & 10.1\% \\
    Health and social assistance & 5.2\% & 6.5\% & 10.6\% \\
    Households as employers & 16.2\% & 0.9\% & 0.0\% \\
    \bottomrule
  \end{tabular}
  \vspace{0.3em}
  \begin{minipage}{0.95\textwidth}
  \footnotesize
  Notes: Statistics use 2025 formal workers and GEIH expansion weights. The first row reports each column's share of all formal workers in the table. The 2--100 column pools all formal workers in firm-size categories from 2--3 through 51--100 workers. Higher education corresponds to workers classified as superior or university educated. Sector rows report the weighted share of each group located in the listed economic activity.
  \end{minipage}
\end{table}


Figure~\ref{fig:minwage-distribution-formality-size} shifts from mean wages to the distribution of hourly wages relative to the statutory hourly minimum in 2025. The vertical line marks one minimum wage, and the labels on the right report each firm-size category's employment share within the corresponding formality group. Among formal workers, the categories from 4--5 to 51--100 workers are visibly concentrated around one minimum wage. Appendix Table~\ref{tab:minimum-wage-bunching-2025} quantifies this pattern: between 36.9 and 42.7 percent of formal workers in these categories earn between 0.9 and 1.1 hourly minimum wages. Formal solo workers and formal workers in 101+ firms have wider upper tails, with 75th percentiles of 2.37 and 2.59 hourly minimum wages, respectively. Among informal workers, the distribution is shifted downward: 73.5 percent of solo workers and 65.9 percent of workers in firms with 2--3 workers earn less than 0.9 hourly minimum wages. This evidence helps interpret the flatter formal-sector profile: the minimum wage appears as a visible compression benchmark in formal wage setting, while informal employment contains much more mass below that benchmark.

\begin{figure}[!htbp]
  \centering
  \includegraphics[width=\textwidth]{fig77_\figurelang.png}
  \caption{Distribution of hourly wages relative to the statutory hourly minimum by firm size and formality status, 2025}
  \label{fig:minwage-distribution-formality-size}
\end{figure}

Figure~\ref{fig:income-sex-comparison} shows that the raw firm-size pattern is
similar for men and women. In 2025, among men, solo workers earned 6,816 pesos
per hour, compared with 16,929 pesos in firms with 101 or more workers; among
women, the corresponding means were 6,133 and 16,777 pesos. The large-firm
advantage is therefore very similar for both sexes. The figure also shows that wage
differences between men and women within a given firm-size category are small
relative to the difference between solo work and large-firm employment.

\begin{figure}[!htbp]
  \centering
  \includegraphics[width=0.95\textwidth]{fig70_\figurelang.png}
  \caption{Mean real hourly income by firm size and sex}
  \label{fig:income-sex-comparison}
\end{figure}

Figure~\ref{fig:income-education-comparison} shows that the raw firm-size
premium is visible within every education group. Among workers with preschool
or less, mean hourly income in 101+ firms is 2.1 times solo-worker income. The
corresponding 101+-to-solo ratio is 1.8 among workers with primary education,
1.5 among workers with secondary education, and 1.7 among workers with higher
education. Thus, the large-firm advantage is not only a comparison between more
and less educated workers; it is also visible within each education level.

\begin{figure}[!htbp]
  \centering
  \includegraphics[width=\textwidth]{fig72_\figurelang.png}
  \caption{Mean real hourly income by firm size and education group}
  \label{fig:income-education-comparison}
\end{figure}

Figure~\ref{fig:income-sector-comparison} shows that the firm-size wage premium
is widespread across economic activities. Mean hourly income in firms with 101 or more workers exceeded
solo-worker income in 14 of 16 sectors in 2008 and in all 16 sectors in 2025.
The two 2008 exceptions were finance and insurance and business services, where
solo workers earned slightly more than workers in the largest firms; by 2025,
the large-firm category also paid more in both sectors. The size premium is
especially pronounced in utilities and water, households as employers, mining,
manufacturing, accommodation and food, and transport, while it is much flatter
in information and communication, health, finance, and business services. These
sector panels therefore reinforce the aggregate pattern while also showing that
the size premium varies substantially across activities.

\begin{figure}[p]
  \centering
  \includegraphics[width=\textwidth]{fig71_\figurelang.png}
  \caption{Mean real hourly income by firm size and sector, 2008 and 2025}
  \label{fig:income-sector-comparison}
\end{figure}

Taken together, the descriptive evidence motivates the regression analysis
below. The central pattern is not only that mean hourly income is higher in
larger firm-size categories. The tables and figures also provide suggestive
evidence that a firm-size wage premium appears within each sex, within each
education group, in most economic activities, and even within both formal and
informal employment. At the same time, the descriptive results show strong
sorting by formality, education, sector, and position in the firm-size
distribution. The econometric exercise therefore asks how much of the raw
firm-size premium survives after controlling for observable worker
characteristics, formality, and sector-year conditions.

\section{Methodology}

\subsection{Main Specification}

The main econometric exercise estimates the conditional firm-size wage premium in a log-wage model. The baseline specification is:
\begin{equation}\label{eq:main}
  \ln(w_{ist}) =
  \alpha
  + \sum_{g \neq g_0} \beta_g \mathbf{1}\{Size_{ist}=g\}
  + X_{ist}'\gamma
  + \lambda_{s \times t}
  + \varepsilon_{ist},    
\end{equation}
where \(w_{ist}\) is real hourly labor income for worker \(i\) in sector \(s\) in year \(t\), \(Size_{ist}\) is the harmonized firm-size category, and \(g_0\) is the reference group. We use solo workers as the reference group. This choice makes each coefficient directly interpretable as the conditional wage gap between workers in firm-size group \(g\) and workers at the bottom of the firm-size distribution. 

The vector \(X_{ist}\) includes a female-worker indicator, age, age squared, education dummies, and a formal-employment indicator. The fixed effect \(\lambda_{s \times t}\) absorbs sector-year conditions, allowing each broad sector to have its own wage level in each year. The model is estimated with GEIH expansion weights, and standard errors are clustered by sector.

The coefficients \(\beta_g\) are log-point differences relative to solo workers. For interpretation, we transform them into percentage premiums:
\begin{equation}
  Premium_g = 100 \times \left(\exp(\hat{\beta}_g)-1\right).    
\end{equation}

\subsection{Dynamic Premiums and Trend Test}

The same framework can be extended to study the evolution of the premium over time by interacting firm-size categories with year indicators:
\begin{equation}
  \ln(w_{ist}) =
  \alpha
  + \sum_{\tau}\sum_{g \neq g_0}
    \beta_{g\tau}\mathbf{1}\{Size_{ist}=g\}\mathbf{1}\{Year_{t}=\tau\}
  + X_{ist}'\gamma
  + \lambda_{s \times t}
  + \varepsilon_{ist}.    
\end{equation}

This dynamic specification produces a sequence of firm-size premiums by year and can be used to assess whether the premium has widened, narrowed, or remained stable. We also estimate a more parsimonious trend specification:
\begin{equation}
  \ln(w_{ist}) =
  \alpha
  + \sum_{g \neq g_0} \left(\beta_g
  +  \delta_g
    Trend_t\right)\mathbf{1}\{Size_{ist}=g\}
  + X_{ist}'\gamma
  + \lambda_{s \times t}
  + \varepsilon_{ist},
\end{equation}
where \(Trend_t\) is a linear year trend normalized to zero in 2008. The parameter \(\delta_g\) measures whether the wage premium for firm-size group \(g\) changed over time relative to solo workers. Positive values of \(\delta_g\) imply a steeper association between firm size and wages, while negative values imply a flatter one.


\subsection{Firm-Size Nonlinearities by Formality Status}

The descriptive evidence suggests that the relationship between firm size and
hourly income may differ between formal and informal workers. We therefore
estimate a specification that interacts all firm-size categories with formal
employment:

\begin{equation}\label{eq:formality-interaction}
  \ln(w_{ist}) =
  \alpha
  + \phi Formal_{ist}
  + \sum_{g \neq g_0} \left(\beta_g 
  + \theta_g
     Formal_{ist}\right)\mathbf{1}\{Size_{ist}=g\}
  + X_{ist}'\gamma
  + \lambda_{s \times t}
  + \varepsilon_{ist}.    
\end{equation}

The coefficients \(\beta_g\) measure the firm-size premium among informal workers, relative to informal solo workers. The corresponding premium among formal workers is \(\beta_g+\theta_g\), relative to formal solo workers. The coefficient \(\phi\) measures the formal-informal wage gap among solo workers, while \(\phi+\theta_g\) measures the formal-informal wage gap within firm-size group \(g\).

\subsection{Approximate Firm-Size Elasticities}

The main specifications treat firm size as a set of categories because this is
how firm size is reported in GEIH and because the largest category is top-coded.
This categorical approach is the most transparent way to estimate the Colombian
firm-size wage premium. However, much of the employer-size wage effect
literature reports a continuous elasticity of wages with respect to employer
size. To compare the Colombian evidence with those benchmarks, we also estimate
an approximate log-log specification:
\begin{equation}\label{eq:elasticity}
  \ln(w_{ist}) =
  \alpha
  + \eta \ln(\widetilde{Size}_{ist})
  + X_{ist}'\gamma
  + \lambda_{s \times t}
  + \varepsilon_{ist}.
\end{equation}

The variable \(\widetilde{Size}_{ist}\) assigns representative values to the
harmonized firm-size bins: 1 for solo workers, 2.5 for firms with 2--3 workers,
4.5 for 4--5 workers, 8 for 6--10 workers, 15 for 11--19 workers, 25 for
20--30 workers, 40.5 for 31--50 workers, 75.5 for 51--100 workers, and 150 for
the top-coded 101+ category. The coefficient \(\eta\) is therefore an
approximate elasticity: it measures the log-wage slope associated with assigned
firm size. Because the underlying variable is grouped and top-coded, these
elasticities are used only as benchmarking statistics. The categorical
specifications remain the main empirical estimates.

\section{Results}

\subsection{Main Results}

Table~\ref{tab:firm-size-wage-premium-regressions} reports the regression results from the main specification in Equation~(\ref{eq:main}). Column (1) shows the raw log-wage gap by firm size. The coefficient for firms with 101 or more workers is 1.019 log points, consistent with the large raw wage premium shown in the descriptive statistics. Column (2) adds sector-year fixed effects, reducing the coefficient for 101+ firms to 0.819 log points. Column (3) adds gender, age, age squared, and education controls, lowering the coefficient to 0.598. Column (4) adds formality, reducing the coefficient further to 0.359 log points, or 43.2 percent. Thus, roughly two thirds of the raw 101+ wage gap is absorbed by sector-year conditions and observed differences in worker and job characteristics, including gender, age, education, and formality. Nonetheless, a positive and statistically significant conditional premium remains after these controls.

\begin{table}[htbp]
  \centering
  \caption{Firm-size wage premium regressions: the role of formality}
  \label{tab:firm-size-wage-premium-regressions}
  \small
  \begin{tabular}{lcccc}
    \toprule
    & \multicolumn{4}{c}{Dependent variable: log real hourly labor income} \\
    \cmidrule(lr){2-5}
    & (1) & (2) & (3) & (4) \\
    \midrule
    Firm size: 2-3 & 0.183*** & 0.242*** & 0.209*** & 0.188*** \\
     & (0.034) & (0.030) & (0.027) & (0.029) \\
    Firm size: 4-5 & 0.363*** & 0.408*** & 0.361*** & 0.307*** \\
     & (0.035) & (0.030) & (0.029) & (0.038) \\
    Firm size: 6-10 & 0.500*** & 0.502*** & 0.422*** & 0.317*** \\
     & (0.029) & (0.033) & (0.037) & (0.050) \\
    Firm size: 11-19 & 0.627*** & 0.577*** & 0.458*** & 0.303*** \\
     & (0.034) & (0.040) & (0.042) & (0.051) \\
    Firm size: 20-30 & 0.689*** & 0.604*** & 0.463*** & 0.271*** \\
     & (0.046) & (0.045) & (0.046) & (0.050) \\
    Firm size: 31-50 & 0.760*** & 0.651*** & 0.488*** & 0.272*** \\
     & (0.056) & (0.049) & (0.047) & (0.048) \\
    Firm size: 51-100 & 0.820*** & 0.699*** & 0.519*** & 0.289*** \\
     & (0.063) & (0.052) & (0.048) & (0.047) \\
    Firm size: 101+ & 1.019*** & 0.819*** & 0.598*** & 0.359*** \\
     & (0.102) & (0.067) & (0.061) & (0.059) \\
    \midrule
    Gender, age, and education controls & No & No & Yes & Yes \\
    Formality control & No & No & No & Yes \\
    Sector-year fixed effects & No & Yes & Yes & Yes \\
    Observations & 5,268,514 & 5,268,514 & 5,268,514 & 5,268,514 \\
    $R^2$ & 0.213 & 0.271 & 0.376 & 0.386 \\
    \bottomrule
  \end{tabular}
  \vspace{0.3em}
  \begin{minipage}{0.95\textwidth}
  \footnotesize
  Notes: The omitted category is solo workers. All columns use the same estimation sample and GEIH expansion weights. Standard errors, clustered by sector, are reported in parentheses. Worker controls include a female-worker indicator, age, age squared, and education dummies. Column (4) adds labor formality. Significance levels: * $p<0.10$, ** $p<0.05$, *** $p<0.01$.
  \end{minipage}
\end{table}


This attenuation pattern is consistent with the classic firm-size wage premium literature. \citet{BrownMedoff1989} show that large employers pay more even after accounting for several observable worker and job characteristics, and \citet{troske1999evidence} reaches a similar conclusion using matched worker-establishment data. Evidence from developing economies also points to sizable large-employer wage gaps. \citet{Schaffner1998} documents premiums to employment in larger establishments in Peru, \citet{SoderbomTealWambugu2005} find that the earnings-firm size relationship remains substantial after accounting for individual fixed effects in Kenya and Ghana, and \citet{reed2019large} show that the large-firm wage premium is widespread across developing economies. The Colombian context adds a further margin because firm size and formality are tightly connected. \citet{ElBadaouiStroblWalsh2010} show for Ecuador that what appears to be a formal-sector wage premium can be largely a firm-size wage differential, and Colombian firm-survey evidence emphasizes the role of firm characteristics in wage differentials within the formal labor market \citep{iregui2012diferenciales}. Our estimates fit this broader pattern. The fall in the coefficient after adding formality shows that the raw firm-size gap partly reflects segmentation between formal and informal jobs; the remaining premium shows that the Colombian firm-size wage gap is not only a formal-informal contrast.

Figure~\ref{fig:firm-size-wage-premium} transforms the coefficients of column (4) in Table~\ref{tab:firm-size-wage-premium-regressions} into percentage premiums. Relative to solo workers, the conditional premium is 20.7 percent for firms with 2--3 workers, 36.0 percent for firms with 4--5 workers, 37.3 percent for firms with 6--10 workers, roughly 31--35 percent for intermediate-size categories, and 43.2 percent for firms with 101 or more workers. All confidence intervals exclude zero.

The choice of solo workers as the reference category is a normalization. Appendix Table~\ref{tab:firm-size-wage-premium-regressions-101-reference} reports the same four specifications using workers in firms with 101 or more workers as the omitted category. In column (4), solo workers earn 0.359 log points less than comparable workers in 101+ firms; this is the log difference underlying the 43.2 percent premium of 101+ workers over solo workers. The alternative normalization also shows that the adjusted gap is largest at the bottom of the firm-size distribution: all smaller categories have negative point estimates relative to 101+ workers, but some middle-size categories are not statistically different from the 101+ category in the full specification.

Appendix Table~\ref{tab:firm-size-wage-premium-regressions-11-19-reference} reports a second normalization using firms with 11--19 workers as the omitted category. In column (4), solo workers and workers in firms with 2--3 workers earn less than otherwise comparable workers in the 11--19 category. The 101+ coefficient is positive, at 0.057 log points, but it is not statistically different from zero with standard errors clustered by sector. This middle-category normalization clarifies that the conditional full-sample premium should not be read as a smooth increase at every step of the firm-size distribution. The strongest adjusted contrast is between solo work and very small firms on one side, and larger firm-size categories on the other.

Appendix Table~\ref{tab:firm-size-local-geography-robustness} shows that this result is not driven by broad geographic differences across departments. The 101+ coefficient remains positive and statistically significant after adding department-year fixed effects, after absorbing sector-department-year fixed effects, and after adding occupational-position controls to that local-geography specification. Appendix Tables~\ref{tab:firm-size-local-geography-robustness-101-reference} and \ref{tab:firm-size-local-geography-robustness-11-19-reference} report the same robustness exercise using 101+ workers and 11--19 workers as the omitted category, respectively.

\begin{figure}[htbp]
  \centering
  \includegraphics[width=0.9\textwidth]{fig57_\figurelang.png}
  \caption{Conditional firm-size wage premium relative to solo workers}
  \label{fig:firm-size-wage-premium}
\end{figure}

\subsection{Dynamic Results}

The dynamic exercise asks whether the firm-size wage gap steepened, flattened,
or remained stable over time. A steepening premium means that firm size became
more strongly associated with wages. A flattening premium means that the
association weakened, although it does not identify whether the change came from
faster wage growth in smaller firms or slower wage growth in larger firms.

Figure~\ref{fig:firm-size-wage-premium-over-time} estimates the full specification separately by interacting firm-size categories with year indicators. The conditional premium is positive throughout the sample period in all firm-size categories. The largest firms display the highest premiums in most years, although the confidence intervals are wider for these categories.

\begin{figure}[htbp]
  \centering
  \includegraphics[width=0.95\textwidth]{fig59_\figurelang.png}
  \caption{Conditional firm-size wage premium over time}
  \label{fig:firm-size-wage-premium-over-time}
\end{figure}

Table~\ref{tab:firm-size-premium-trend-test} formally tests whether the premium changed over time using a linear trend interacted with firm-size categories. The estimated trend coefficients are positive for all size categories. The implied increase between 2008 and 2025 ranges from 3.2 percent for firms with 20--30 workers to 7.9 percent for firms with 101 or more workers. However, the two-sided tests do not reject a zero trend at the 5 percent level. The point estimates therefore suggest widening rather than narrowing, but they do not provide statistically conclusive evidence of a trend.

Appendix Table~\ref{tab:firm-size-premium-trend-test-101-reference} reports the same trend test with 101+ workers as the reference group. In that normalization, a negative coefficient means that the listed category lost ground relative to 101+ workers. The table gives the same substantive message from the top of the firm-size distribution: the point estimates suggest that several categories lost ground relative to 101+ workers, but the evidence is not strong enough to conclude that the premium systematically widened across the full distribution.

Appendix Table~\ref{tab:firm-size-premium-trend-test-11-19-reference} uses the 11--19 category as the reference group. Relative to this middle category, the implied 2008--2025 change is 3.0 percent for 101+ workers and is not statistically significant. The mid-reference test therefore reinforces the same conclusion: the point estimates suggest some widening at the top, but the evidence does not support a statistically conclusive trend.

The evidence on the evolution of the large-firm wage premium over time is mixed. For the United States, \citet{BloomGuvenenSmithSongVonWachter2018} document a steady decline in the large-firm wage premium. \citet{reed2019large} shows that large-firm wage premia are also common in developing economies, but tend to decline with national income and over time. Cross-country evidence points to heterogeneity rather than a single trajectory: \citet{ColonnelliTagWebbWolter2018} find a large-firm wage premium in Brazil, Germany, Sweden, and the United Kingdom, but no common time path across countries. Evidence from Brazil also shows that firm-related wage differences can compress when institutions and labor-market conditions raise wages at the lower end of the distribution \citep{AlvarezBenguriaEngbomMoser2018,EngbomMoser2022}. Against this background, Colombia looks less like a case of clear convergence or clear divergence and more like a case of persistence: firm scale remains strongly associated with hourly wages throughout the period, the point estimates suggest some widening, but the two-sided tests do not support the stronger claim that the premium systematically intensified.

\begin{table}[htbp]
  \centering
  \caption{Testing whether the firm-size wage premium changed over time}
  \label{tab:firm-size-premium-trend-test}
  \small
  \begin{tabular}{lcccc}
    \toprule
    Firm size & Annual trend & S.E. & $p$-value & Implied 2008--2025 change \\
    \midrule
    2-3 & 0.0033 & (0.0028) & 0.262 & 5.7\% \\
    4-5 & 0.0032 & (0.0031) & 0.314 & 5.6\% \\
    6-10 & 0.0021 & (0.0034) & 0.536 & 3.7\% \\
    11-19 & 0.0027 & (0.0039) & 0.494 & 4.7\% \\
    20-30 & 0.0018 & (0.0041) & 0.661 & 3.2\% \\
    31-50 & 0.0027 & (0.0044) & 0.542 & 4.7\% \\
    51-100 & 0.0026 & (0.0042) & 0.549 & 4.5\% \\
    101+ & 0.0045 & (0.0043) & 0.314 & 7.9\% \\
    \bottomrule
  \end{tabular}
  \vspace{0.3em}
  \begin{minipage}{0.95\textwidth}
  \footnotesize
  Notes: The annual trend is the coefficient on the interaction between firm-size category and a linear year trend, with solo workers as the omitted category. The specification controls for gender, age, age squared, education, formality, and sector-year fixed effects, and uses GEIH expansion weights. Standard errors are clustered by sector. The $p$-value corresponds to the two-sided test that the trend differs from zero. The final column reports $100\times[\exp(17\hat{\delta})-1]$, the implied change in the firm-size wage ratio between 2008 and 2025.
  \end{minipage}
\end{table}


\subsection{Formality Nonlinearities}

Table~\ref{tab:formality-interaction-coefficients} reports the coefficients behind Equation~(\ref{eq:formality-interaction}). The estimates of \(\beta_g\), which measure the firm-size premium among informal workers, are positive and statistically significant in every firm-size category. The estimates of \(\theta_g\) are negative in every category, showing that the formal-worker profile is flatter than the informal-worker profile. The implied formal-worker premium \(\beta_g+\theta_g\) is small and not statistically different from zero at the 5 percent level when formal solo workers are the reference category. This pattern shows a robust size premium among informal workers but a much flatter profile among formal workers.

\begin{table}[htbp]
  \centering
  \caption{Firm-size coefficients by formality status}
  \label{tab:formality-interaction-coefficients}
  \small
  \begin{tabular}{lccc}
    \toprule
    Firm size & $\hat{\beta}_g$ & $\hat{\theta}_g$ & $\hat{\beta}_g+\hat{\theta}_g$ \\
    \midrule
    2-3 & 0.197*** & -0.176*** & 0.021 \\
     & (0.026) & (0.062) & (0.067) \\
    4-5 & 0.343*** & -0.318*** & 0.024 \\
     & (0.027) & (0.085) & (0.089) \\
    6-10 & 0.392*** & -0.369*** & 0.023 \\
     & (0.032) & (0.081) & (0.087) \\
    11-19 & 0.409*** & -0.374*** & 0.035 \\
     & (0.037) & (0.076) & (0.080) \\
    20-30 & 0.402*** & -0.379*** & 0.023 \\
     & (0.035) & (0.072) & (0.079) \\
    31-50 & 0.401*** & -0.359*** & 0.041 \\
     & (0.038) & (0.067) & (0.076) \\
    51-100 & 0.410*** & -0.342*** & 0.068 \\
     & (0.044) & (0.067) & (0.075) \\
    101+ & 0.336*** & -0.183** & 0.153* \\
     & (0.046) & (0.083) & (0.088) \\
    \bottomrule
  \end{tabular}
  \vspace{0.3em}
  \begin{minipage}{0.92\textwidth}
  \footnotesize
  Notes: Estimates correspond to Equation~(\ref{eq:formality-interaction}). The omitted category is solo workers within each formality group. $\hat{\beta}_g$ is the firm-size coefficient among informal workers. $\hat{\theta}_g$ is the additional coefficient for formal workers in the same firm-size category. $\hat{\beta}_g+\hat{\theta}_g$ is the implied firm-size coefficient among formal workers. All coefficients are in log points. Standard errors, clustered by sector, are reported in parentheses. The specification controls for gender, age, age squared, education, and sector-year fixed effects. Significance levels: * $p<0.10$, ** $p<0.05$, *** $p<0.01$.
  \end{minipage}
\end{table}


Figure~\ref{fig:firm-size-formality-premium} translates these estimates into percentage premiums relative to solo workers within each formality group. Among informal workers, the conditional premium is positive and statistically significant in every firm-size category: it is 21.7 percent in firms with 2--3 workers, rises to about 48--51 percent across most categories from 6 to 100 workers, and remains 40.0 percent in firms with 101 or more workers. Among formal workers, the within-formality profile is much flatter and less precisely estimated: the estimated premium ranges from 2.1 to 7.0 percent for most categories and is 16.5 percent for firms with 101 or more workers, but the confidence intervals include zero throughout when each category is compared with formal solo workers.

\begin{figure}[htbp]
  \centering
  \includegraphics[width=0.95\textwidth]{fig74_\figurelang.png}
  \caption{Conditional firm-size wage premium by formality status}
  \label{fig:firm-size-formality-premium}
\end{figure}

This does not mean, however, that there is no evidence of conditional premiums within the formal sector. The U-shape in Figure~\ref{fig:income-formality-comparison} points to a different comparison: formal workers in 101+ firms versus formal workers in intermediate-size firms. Direct contrasts from the interacted model show that formal workers in 101+ firms earn between 8.9 and 14.2 percent more than formal workers in every smaller non-solo category, and all of these non-solo contrasts are statistically significant, as reported in Table~\ref{tab:formal-101-contrasts}. Thus, the estimates show a large-firm premium relative to the middle of the formal firm-size distribution. The informal sector, by contrast, displays a robust conditional firm-size premium relative to solo work.

\begin{table}[htbp]
  \centering
  \caption{Formal-worker large-firm wage premium relative to other formal firm-size categories}
  \label{tab:formal-101-contrasts}
  \small
  \begin{tabular}{lcc}
    \toprule
    Reference category & Premium & $p$-value \\
    \midrule
    Solo workers & 16.5\% & 0.082 \\
    2-3 & 14.2\% & $<0.001$ \\
    4-5 & 13.8\% & $<0.001$ \\
    6-10 & 13.9\% & $<0.001$ \\
    11-19 & 12.6\% & 0.004 \\
    20-30 & 13.9\% & 0.001 \\
    31-50 & 11.8\% & 0.002 \\
    51-100 & 8.9\% & 0.004 \\
    \bottomrule
  \end{tabular}
  \vspace{0.3em}
  \begin{minipage}{0.92\textwidth}
  \footnotesize
  Notes: Each row compares formal workers in firms with 101 or more workers with formal workers in the reference firm-size category. Estimates come from the firm-size-by-formality specification with gender, age, age squared, education, and sector-year fixed effects. Premiums are computed as $100\times[\exp(\hat{\beta})-1]$ from the relevant linear contrast.
  \end{minipage}
\end{table}


Appendix Table~\ref{tab:formality-interaction-coefficients-101-reference} reports the interacted specification using 101+ workers as the reference category within each formality group. This alternative normalization clarifies the shape of the within-formality profiles. Among formal workers, the categories from 2--3 to 51--100 workers are below formal 101+ workers, which is the comparison reported in Table~\ref{tab:formal-101-contrasts}. Among informal workers, the lowest categories are below informal 101+ workers, but several middle-size categories are close to or above the 101+ coefficient. The informal premium should therefore be read as a large wage gap between solo and very small informal units and the rest of the informal firm-size distribution, not as a strictly monotonic increase up to 101+ firms.

Appendix Table~\ref{tab:formality-interaction-coefficients-11-19-reference} uses the 11--19 category as the reference group within each formality group. Among informal workers, solo work and firms with 2--3 workers are clearly below the 11--19 category, while the categories above 11--19 are economically close to that reference. Among formal workers, the 101+ category is 0.119 log points above the 11--19 category and statistically significant. The middle-reference table therefore sharpens the main formality result: the informal profile rises sharply out of solo work and the smallest firms, while the formal profile is flatter except for the advantage of 101+ firms over the middle formal categories.

Table~\ref{tab:formal-wage-gap-by-size} reports the other side of the same specification: the conditional wage gap between formal and informal workers within each firm-size category. Formal workers earn more than informal workers in every category, and the differences are statistically significant throughout. The formal premium is 68.8 percent among solo workers, falls to roughly 16--23 percent through most intermediate firm-size categories, and rises again to 40.6 percent among workers in firms with 101 or more employees. Thus, the flatter formal-worker firm-size profile does not mean that formality is unimportant. Rather, formal employment shifts the wage level upward within each firm-size category, while the firm-size premium itself is much steeper among informal workers.

\begin{table}[htbp]
  \centering
  \caption{Conditional formal-informal wage gap by firm size}
  \label{tab:formal-wage-gap-by-size}
  \small
  \begin{tabular}{lccc}
    \toprule
    Firm size & Log gap & Premium & $p$-value \\
    \midrule
    Solo & 0.524*** & 68.8\% & $<0.001$ \\
     & (0.063) &  &  \\
    2-3 & 0.348*** & 41.6\% & $<0.001$ \\
     & (0.018) &  &  \\
    4-5 & 0.205*** & 22.8\% & $<0.001$ \\
     & (0.036) &  &  \\
    6-10 & 0.155*** & 16.7\% & $<0.001$ \\
     & (0.031) &  &  \\
    11-19 & 0.149*** & 16.1\% & $<0.001$ \\
     & (0.024) &  &  \\
    20-30 & 0.145*** & 15.6\% & $<0.001$ \\
     & (0.018) &  &  \\
    31-50 & 0.165*** & 17.9\% & $<0.001$ \\
     & (0.014) &  &  \\
    51-100 & 0.181*** & 19.9\% & $<0.001$ \\
     & (0.019) &  &  \\
    101+ & 0.341*** & 40.6\% & $<0.001$ \\
     & (0.043) &  &  \\
    \bottomrule
  \end{tabular}
  \vspace{0.3em}
  \begin{minipage}{0.92\textwidth}
  \footnotesize
  Notes: Each row reports the estimated wage gap between formal and informal workers within the same firm-size category from Equation~(\ref{eq:formality-interaction}). For solo workers, the gap is $\hat{\phi}$; for all other categories, it is $\hat{\phi}+\hat{\theta}_g$. Premiums are computed as $100\times[\exp(\widehat{gap})-1]$. Standard errors, clustered by sector, are reported in parentheses below the log gap. The specification controls for gender, age, age squared, education, and sector-year fixed effects. Significance levels: * $p<0.10$, ** $p<0.05$, *** $p<0.01$.
  \end{minipage}
\end{table}


Figure~\ref{fig:adjusted-formality-endpoints} asks whether the adjusted formality-specific profiles changed between 2008 and 2025. The figure estimates the firm-size-by-formality specification separately for each endpoint year, controlling for gender, age, age squared, education, and sector fixed effects, and plots the implied premium relative to solo workers within the same formality-year group. Among formal workers, the adjusted profile is flat in both years: in 2025, the estimates for all categories below 101+ range from -3.7 to 0.6 percent, and the 101+ estimate is 10.5 percent with a confidence interval that includes zero. Among informal workers, by contrast, the adjusted size profile is steep in both years and higher at the 2025 endpoint: in 2025, the premium rises from 24.3 percent in firms with 2--3 workers to 53.6 percent in firms with 101 or more workers. The figure therefore suggests that the strong conditional size premium within informal employment persists across the endpoints, while the formal-sector profile remains flat.

\begin{figure}[htbp]
  \centering
  \includegraphics[width=0.95\textwidth]{fig76_\figurelang.png}
  \caption{Adjusted firm-size wage premium by formality status, 2008 and 2025}
  \label{fig:adjusted-formality-endpoints}
\end{figure}



Figure~\ref{fig:large-firm-formality-premium-over-time} examines whether the
pooled formal-worker result masks time variation. The figure plots the adjusted
premium for firms with 101 or more workers relative to solo workers, separately
within formal and informal employment in each year. The evidence does not show a
reappearance of a statistically significant large-firm premium among formal
workers in the most recent years. The formal 101+ premium is positive and
significant in several earlier years, especially 2009--2013 and 2017, but after
2021 the point estimates are smaller and the confidence intervals include zero;
in 2025, the estimate is 10.5 percent with a confidence interval from -8.2 to
33.0 percent. By contrast, the informal 101+ premium is positive and significant
throughout the period, rising from 39.4 percent in 2008 to 53.6 percent in
2025.

\begin{figure}[htbp]
  \centering
  \includegraphics[width=0.95\textwidth]{fig75_\figurelang.png}
  \caption{Adjusted large-firm wage premium by formality status over time; 2020 is omitted}
  \label{fig:large-firm-formality-premium-over-time}
\end{figure}

This formal-informal contrast connects directly with the literature on formality and firm size. Formality can mediate the firm-size wage premium because larger firms are more visible to the state, more likely to comply with labor regulation, and more likely to offer jobs covered by social protection. Colombian evidence already points to persistent formal-informal wage gaps and to productivity differences associated with registration and regulatory compliance among microfirms \citep{DazaGamboa2013,GutierrezRodriguezLesmes2023}. In that sense, part of the raw premium may reflect the fact that large-firm jobs are disproportionately formal jobs. From the complementary angle, \citet{ElBadaouiStroblWalsh2010} show for Ecuador that what appears to be a formal-sector wage premium can be largely explained by firm size. Formality can also moderate the firm-size premium: the wage gap by firm size need not be the same within formal and informal employment. \citet{BalkanTumen2016} show for Turkey that the firm-size wage gap is more pronounced among informal jobs than among formal jobs. Our results are close to this second pattern. Adding formality in Table~\ref{tab:firm-size-wage-premium-regressions} reduces the overall 101+ coefficient substantially, which shows that formality mediates part of the aggregate size gap. Figures~\ref{fig:firm-size-formality-premium} and \ref{fig:adjusted-formality-endpoints} show the complementary result: the conditional size premium remains especially strong among informal workers, while the formal-worker profile is much flatter when each category is compared with formal solo workers.

\subsection{Benchmarking the Firm-Size Wage Elasticity}
\label{subsec:elasticity-benchmark}

Table~\ref{tab:firm-size-elasticity-benchmark} reports the approximate
elasticity estimates from Equation~(\ref{eq:elasticity}) and compares them with
the two most directly comparable benchmarks from the meta-analysis of
\citet{DiegmannMullerSchoefer2026}: specifications without worker-heterogeneity
controls and specifications with observable worker controls. The comparison
should be read as a magnitude check rather than as a direct meta-analytic test:
the Colombian estimates are based on worker-reported firm-size bins, while the
studies summarized by \citet{DiegmannMullerSchoefer2026} estimate continuous
employer-size elasticities with different data structures and controls.

\begin{table}[htbp]
  \centering
  \caption{Benchmarking Colombia's approximate firm-size wage elasticity}
  \label{tab:firm-size-elasticity-benchmark}
  \small
  \begin{tabular}{p{0.30\textwidth}p{0.40\textwidth}cc}
    \toprule
    Source & Specification & Elasticity & S.E. \\
    \midrule
    Diegmann et al. (2026) & No worker-heterogeneity controls & 0.062 & n.a. \\
    Diegmann et al. (2026) & Observable worker controls & 0.036 & n.a. \\
    \midrule
    Colombia (GEIH) & No controls & 0.203*** & (0.019) \\
    Colombia (GEIH) & Full sample & 0.071*** & (0.013) \\
    Colombia (GEIH) & Formal workers & 0.026** & (0.011) \\
    Colombia (GEIH) & Informal workers & 0.130*** & (0.015) \\
    \bottomrule
  \end{tabular}
  \vspace{0.3em}
  \begin{minipage}{0.95\textwidth}
  \footnotesize
  Notes: Diegmann et al. (2026) benchmarks correspond to the mean elasticities reported in their Table 1 for specifications without worker-heterogeneity controls and with observable worker controls. Their worker fixed-effect and AKM employer-effect benchmarks are not reported because GEIH does not follow workers across employers or identify firms. Colombian elasticities are approximate log-log slopes estimated with GEIH expansion weights. The full-sample Colombian specification includes sector-year fixed effects, worker controls, and formality; the formal- and informal-worker specifications include sector-year fixed effects and worker controls. Because GEIH reports firm size in bins, the estimates assign representative values of 1, 2.5, 4.5, 8, 15, 25, 40.5, 75.5, and 150 workers to the harmonized categories from solo work to 101+ workers. Standard errors for the Colombian estimates are clustered by sector. Significance levels: * $p<0.10$, ** $p<0.05$, *** $p<0.01$.
  \end{minipage}
\end{table}


The Colombian raw binned elasticity is 0.203, far above the mean no-controls
elasticity of 0.062 in the literature summarized by
\citet{DiegmannMullerSchoefer2026}. After sector-year effects, worker controls,
and formality are added, the full-sample Colombian elasticity falls to 0.071,
but remains above their observable-controls benchmark of 0.036. The formality
split helps interpret why the Colombian full-sample elasticity is high relative
to that benchmark. The difference is not driven equally by formal and informal
employment. Among formal workers, the estimated elasticity is 0.026, close to
the benchmark. Among informal workers, it is 0.130, more than three times the
benchmark. Thus, the comparison with the meta-analysis clarifies that
Colombia's high conditional elasticity is mainly an informal-sector phenomenon.

\section{Discussion}

The results link firm size to three related outcomes: wage inequality,
labor-market segmentation, and possible allocative frictions. The inequality
result is direct: workers in larger firm-size categories earn more than workers
in solo work and micro-firms, both in raw averages and after conditioning on
observed characteristics. The segmentation result comes from the role of
formality. Adding formality to the main regression absorbs part of the raw
large-firm premium, and formal workers earn more than informal workers within
every firm-size category. At the same time, the premium does not disappear after
controlling for formality, and the size profile differs sharply between formal
and informal employment.

The formality results narrow the policy-relevant interpretation. Formal employment raises wage levels within every firm-size category, but the conditional size premium is steepest among informal workers. Within formal employment, the premium relative to formal solo workers is small and not statistically significant, although workers in 101+ firms earn more than workers in intermediate-size formal firms. If the conditional wage gaps partly reflect productivity-related frictions rather than only unobserved sorting, the evidence points to three margins for possible productivity gains: moving workers and jobs from informality to formality, scaling informal micro-units into larger and more organized firms, and, more selectively, growth or movement within the formal sector from intermediate-size firms toward larger employers.

These interpretations must remain cautious, though. A wage premium that survives observable controls does not identify a single mechanism as an explanation. As \citet{BrownMedoff1989} and \citet{OiIdson1999} emphasize, the remaining gap may reflect unobserved worker quality, firm productivity, rent sharing, efficiency wages, or match quality. The developing-country evidence cautions against reducing the premium to sorting alone: \citet{SoderbomTealWambugu2005} find that the firm-size effect remains substantial in Kenya and Ghana after controlling for individual fixed effects. Matched employer-employee studies also show that employer-associated characteristics matter for wage variation \citep{AbowdKramarzMargolis1999,CardHeiningKline2013,Song2019}, and related evidence from Brazil reaches a similar conclusion for Latin America \citep{AlvarezBenguriaEngbomMoser2018,GerardLagosSeverniniCard2021, EngbomMoser2022}. Recent German evidence sharpens this caution: \citet{DiegmannMullerSchoefer2026} find that the employer-size wage effect disappears when establishment size is compared within the same multi-establishment firm, suggesting that the premium may reflect employer-wide wage policies, rents, or other firm attributes rather than establishment size per se. The elasticity benchmark in Section~\ref{subsec:elasticity-benchmark} is consistent with a large Colombian size-wage association, especially among informal workers, but GEIH does not identify employers or follow workers across firms. The paper therefore cannot separate worker sorting from other competing mechanisms, so the estimated premiums and elasticities should be interpreted as diagnostic magnitudes rather than causal effects of firm size.

The evidence is also not a direct test of misallocation. Misallocation in the strict sense requires comparing marginal products across firms or sectors, as in \citet{hsieh2009misallocation}. The link with allocative efficiency comes from the marginal-product condition: in an efficient allocation, labor should move across firms until the marginal product of labor is equalized across uses. In the competitive benchmark, wages reflect the value of marginal product, so persistent wage differences across firm sizes are informative about possible allocative frictions. If larger firms pay more than smaller firms for workers with the same observable characteristics, this may indicate that labor is more valuable in larger firms or that barriers prevent workers from moving toward higher-productivity matches. The evidence is only a marker, not proof, because wages may also reflect unobserved worker quality or other unobserved features of the employment relationship. Testing misallocation directly would require firm-level data on output, capital, and employment, which would allow us to estimate marginal products rather than infer them indirectly from wages.

The policy implication is therefore not that policy should favor large firms as such or push workers mechanically into large firms. The relevant objective is to reduce barriers that keep productive firms small, informal, or disconnected from larger markets, while also improving workers' access to better firm matches. The formality results sharpen this point. Because formal employment raises wage levels within every firm-size category, moving jobs from informality to formality remains a policy priority. Because the firm-size premium is steepest among informal workers, a second policy priority is helping productive informal micro-units become larger, more organized, and more connected to markets. This support should be designed as a path toward formalization, not as a subsidy for remaining informal; otherwise, it could weaken incentives for formal firms to remain formal. The evidence gives less support to a policy focus on scaling up formal firms: within formal employment, the premium relative to formal solo workers is small, although workers in 101+ firms earn more than workers in intermediate-size formal firms. This pattern is consistent with a segmented labor market in which the strongest wage-based signs of allocative frictions are concentrated in informal employment. The largest efficiency gains therefore appear to lie in formalization and in the growth of productive informal firms, rather than in firm growth within the formal sector.

\section{Conclusion}

This paper documents a clear firm-size wage premium in Colombia using harmonized GEIH microdata for 2008--2019 and 2021--2025. The descriptive evidence shows that real hourly wages rise sharply with firm size, while employment remains concentrated in solo work and micro-firms. Solo workers and workers in firms with up to ten workers account for about two thirds of employment, and the middle of the firm-size distribution remains thin.

The regression results show that a substantial conditional premium remains after controlling for gender, age, age squared, education, formality, and sector-year fixed effects. In the full specification, workers in firms with 101 or more employees earn about 43 percent more than otherwise comparable solo workers. The premium is positive throughout the sample period. The point estimates suggest an increase between 2008 and 2025, but the estimated change is not statistically significant.

The formal-informal results sharpen the main finding. Formal workers earn more than informal workers within every firm-size category. At the same time, the firm-size premium is much steeper among informal workers than among formal workers. Within formal employment, the premium relative to formal solo workers is small, although workers in 101+ firms earn more than workers in intermediate-size formal firms. Taken together, these results show that the Colombian firm-size wage premium is
not only a formal-informal contrast. It is a persistent wage gap associated with firm scale, especially within informal employment, and it remains after conditioning on the main observable worker and job characteristics available in
the GEIH.

%The approximate elasticity estimates reinforce this conclusion. The full-sample conditional elasticity is 0.071, above the observable-controls benchmark in the meta-analysis of \citet{DiegmannMullerSchoefer2026}, but the split by formality shows that this relatively high elasticity is mainly an informal-sector phenomenon: the estimate is 0.026 among formal workers and 0.130 among informal workers.


\appendix
\section{Additional Regression Table}
\label{app:additional-regression-table}

This appendix reports the robustness checks referenced in Section~5. Table~\ref{tab:firm-size-local-geography-robustness} first adds department-level fixed effects. In Colombia, departments are first-level territorial units, with Bogot\'a treated as a department-equivalent Capital District. The most demanding geography specification absorbs sector-department-year fixed effects, so the firm-size coefficients are identified within the same sector, department, and year. The final column adds occupational-position controls to this specification. The large-firm coefficient remains positive and statistically significant: workers in firms with 101 or more employees earn about 40.2 percent more than comparable solo workers after conditioning on sector-department-year cells and occupational position.

\begin{table}[htbp]
  \centering
  \caption{Firm-size wage premium robustness checks}
  \label{tab:firm-size-local-geography-robustness}
  \small
  \begin{tabular}{lcccc}
    \toprule
    & \multicolumn{4}{c}{Dependent variable: log real hourly labor income} \\
    \cmidrule(lr){2-5}
    & (1) & (2) & (3) & (4) \\
    \midrule
    Firm size: 2-3 & 0.188*** & 0.168*** & 0.167*** & 0.102*** \\
     & (0.011) & (0.011) & (0.011) & (0.013) \\
    Firm size: 4-5 & 0.307*** & 0.277*** & 0.273*** & 0.226*** \\
     & (0.014) & (0.014) & (0.014) & (0.015) \\
    Firm size: 6-10 & 0.316*** & 0.282*** & 0.278*** & 0.265*** \\
     & (0.018) & (0.017) & (0.017) & (0.016) \\
    Firm size: 11-19 & 0.302*** & 0.272*** & 0.267*** & 0.272*** \\
     & (0.018) & (0.017) & (0.017) & (0.015) \\
    Firm size: 20-30 & 0.271*** & 0.241*** & 0.237*** & 0.258*** \\
     & (0.019) & (0.017) & (0.017) & (0.015) \\
    Firm size: 31-50 & 0.271*** & 0.245*** & 0.240*** & 0.264*** \\
     & (0.018) & (0.017) & (0.017) & (0.015) \\
    Firm size: 51-100 & 0.289*** & 0.262*** & 0.257*** & 0.282*** \\
     & (0.018) & (0.017) & (0.017) & (0.015) \\
    Firm size: 101+ & 0.359*** & 0.338*** & 0.325*** & 0.338*** \\
     & (0.023) & (0.024) & (0.024) & (0.022) \\
    \midrule
    Worker and formality controls & Yes & Yes & Yes & Yes \\
    Occupational-position controls & No & No & No & Yes \\
    Fixed effects & $s\times t$ & $s\times t$, $d\times t$ & $s\times d\times t$ & $s\times d\times t$ \\
    Observations & 5,268,514 & 5,268,514 & 5,268,375 & 5,268,375 \\
    $R^2$ & 0.386 & 0.404 & 0.414 & 0.429 \\
    \bottomrule
  \end{tabular}
  \vspace{0.3em}
  \begin{minipage}{0.95\textwidth}
  \footnotesize
  Notes: The omitted category is solo workers. The sample is restricted to observations with nonmissing department codes. Worker controls include a female-worker indicator, age, age squared, education dummies, and labor formality. The local geography $d$ is the department, treating Bogot\'a as a department-equivalent Capital District. Column (3) absorbs sector-department-year fixed effects, so identification comes from wage differences across firm-size categories within the same sector, department, and year. Column (4) adds occupational-position controls to this most stringent local-geography specification. Standard errors are clustered by sector-department cells. Significance levels: * $p<0.10$, ** $p<0.05$, *** $p<0.01$.
  \end{minipage}
\end{table}



\section*{Data and Code Availability}

Replication code and materials will be made available in a public repository. Raw GEIH microdata are not redistributed and should be obtained directly from DANE.

\bibliographystyle{apalike}
\bibliography{Paper/references}

\end{document}
